% 分类目录：ml
\documentclass[12pt, a4paper]{article}
\usepackage[margin=2.5cm]{geometry}
\usepackage{ctex}
\usepackage{amsmath, amssymb}
\usepackage{listings}
\usepackage{xcolor}
\usepackage{tabularx}
\usepackage{hyperref}

\lstset{
    language=Python,
    basicstyle=\ttfamily\small,
    keywordstyle=\color{blue},
    commentstyle=\color{green!50!black},
    stringstyle=\color{red},
    showstringspaces=false,
    breaklines=true,
    numbers=left,
    numberstyle=\tiny\color{gray},
    frame=single
}

\title{知识点：逻辑回归 (Logistic Regression)}
\author{}
\date{}

\begin{document}
\maketitle

\section{核心定义}
\textbf{中文定义}：逻辑回归（Logistic Regression）是一种广义线性模型，虽然名字含有"回归"，但它主要用于\textbf{分类}任务。其核心思想是通过 Sigmoid 函数将线性组合的输出映射到 $(0, 1)$ 区间，表示样本属于某一类的概率。

\textbf{英文定义}：Logistic Regression is a generalized linear model used primarily for classification. It models the probability of the positive class by applying the logistic (sigmoid) function to a linear combination of input features.

\textbf{历史地位}：逻辑回归是机器学习中最经典、最基础的分类算法之一，也是深度学习中神经元和交叉熵损失函数的直接来源。理解逻辑回归是理解神经网络的重要前提。

\section{数学建模：从线性回归到逻辑回归}

\subsection{为什么不能直接用线性回归做分类？}
线性回归的输出 $\hat{y} = w^\top x + b$ 范围是 $(-\infty, +\infty)$，无法直接解释为概率。若强行用阈值划分，当数据分布变化时，决策边界会剧烈偏移，鲁棒性极差。

\subsection{对数几率视角：Sigmoid 从何而来？}
逻辑回归的名字来源于 \textbf{Logit}（对数几率），其推导过程比直接"选一个 Sigmoid"更加自然和深刻。

\textbf{Step 1: 定义几率 (Odds)}。
对于一个事件发生的概率 $p$，\textbf{几率}定义为该事件发生与不发生的概率之比：
$$ \text{Odds} = \frac{p}{1 - p} $$
几率的范围是 $(0, +\infty)$。例如，$p=0.8$ 时 Odds $= 4$，表示"发生的可能性是不发生的 4 倍"。

\textbf{Step 2: 取对数得到 Logit}。
几率的范围是 $(0, +\infty)$，仍然不对称。对其取自然对数：
$$ \text{logit}(p) = \ln \frac{p}{1 - p} $$
此时 logit 的值域变为 $(-\infty, +\infty)$，与线性模型的输出空间\textbf{完美匹配}！

\textbf{Step 3: 令 logit 等于线性模型}。
这是逻辑回归最核心的建模假设——我们假设\textbf{对数几率}是自变量的线性函数：
$$ \ln \frac{p}{1-p} = w^\top x + b $$
这个假设非常自然：线性模型擅长建模 $(-\infty, +\infty)$ 范围的量，而 logit 恰好就是这样的量。

\textbf{Step 4: 反解出 Sigmoid}。
从上式解出 $p$：
\begin{align}
\frac{p}{1-p} &= e^{w^\top x + b} \\
p &= e^{w^\top x + b} (1-p) \\
p (1 + e^{w^\top x + b}) &= e^{w^\top x + b} \\
p &= \frac{e^{w^\top x + b}}{1 + e^{w^\top x + b}} = \frac{1}{1 + e^{-(w^\top x + b)}} = \sigma(w^\top x + b)
\end{align}
\textbf{Sigmoid 函数就这样自然地"涌现"了}——它不是人为选择的，而是从"对数几率是线性函数"这一假设中数学推导出来的。

\subsection{Sigmoid 函数的关键性质}
逻辑回归引入 Sigmoid（Logistic）函数：
$$ \sigma(z) = \frac{1}{1 + e^{-z}} $$
\textbf{变量含义}：
\begin{itemize}
    \item $z = w^\top x + b$：线性组合的输出，称为 \textbf{logit}。
    \item $w \in \mathbb{R}^d$：权重向量，$d$ 为特征维度。
    \item $b \in \mathbb{R}$：偏置项。
    \item $\sigma(z) \in (0, 1)$：输出的概率值。
\end{itemize}
\textbf{关键性质}：
\begin{enumerate}
    \item $\sigma(0) = 0.5$（决策边界）。
    \item $\sigma(-z) = 1 - \sigma(z)$（对称性）。
    \item 导数形式优雅：$\sigma'(z) = \sigma(z)(1 - \sigma(z))$，这使得梯度计算非常简洁。
\end{enumerate}

\subsection{概率模型}
逻辑回归假设：
$$ P(y=1|x) = \sigma(w^\top x + b), \quad P(y=0|x) = 1 - \sigma(w^\top x + b) $$
可以统一写为：
$$ P(y|x) = \hat{y}^y (1 - \hat{y})^{1-y}, \quad \hat{y} = \sigma(w^\top x + b) $$

\section{损失函数推导：为什么用交叉熵？}

\subsection{极大似然估计 (MLE)}
给定 $N$ 个独立同分布的训练样本 $\{(x_i, y_i)\}$，似然函数为：
$$ L(w, b) = \prod_{i=1}^N \hat{y}_i^{y_i} (1 - \hat{y}_i)^{1 - y_i} $$
取负对数得到\textbf{交叉熵损失函数}（Binary Cross-Entropy, BCE）：
$$ \mathcal{L}(w, b) = -\frac{1}{N} \sum_{i=1}^N \left[ y_i \log \hat{y}_i + (1 - y_i) \log(1 - \hat{y}_i) \right] $$
\textbf{为什么不用 MSE？}
\begin{itemize}
    \item MSE 与 Sigmoid 结合后损失函数是\textbf{非凸}的，存在大量局部最优。
    \item 交叉熵损失函数关于 $w$ 是\textbf{凸函数}，保证全局最优解唯一。
    \item 交叉熵的梯度形式更简洁，训练效率更高。
\end{itemize}

\subsection{梯度推导}
对权重 $w_j$ 求偏导：
$$ \frac{\partial \mathcal{L}}{\partial w_j} = \frac{1}{N} \sum_{i=1}^N (\hat{y}_i - y_i) \cdot x_{ij} $$
\textbf{推导要点}：
\begin{enumerate}
    \item 利用 $\sigma'(z) = \sigma(z)(1 - \sigma(z))$。
    \item 对数求导后与 Sigmoid 导数相消，最终形式极其简洁。
    \item 更新规则：$w_j \leftarrow w_j - \alpha \cdot \frac{\partial \mathcal{L}}{\partial w_j}$，其中 $\alpha$ 是学习率。
\end{enumerate}
这个梯度形式与线性回归的梯度形式相同（$\hat{y}_i - y_i$ 乘以输入），是逻辑回归数学优雅性的体现。

\section{决策边界}
逻辑回归的决策边界由 $w^\top x + b = 0$ 定义，这是一个\textbf{线性超平面}。
\begin{itemize}
    \item 在二维空间中，决策边界是一条直线。
    \item 在高维空间中，决策边界是一个超平面。
    \item 这意味着逻辑回归本质上是一个\textbf{线性分类器}，无法处理非线性可分的数据（如 XOR 问题）。
\end{itemize}
\textbf{解决非线性}：可以通过手动构造多项式特征或使用核技巧（Kernel Trick）来扩展。

\section{正则化}
为防止过拟合，通常在损失函数中加入正则项：
\begin{itemize}
    \item \textbf{L2 正则化}（Ridge）：$\mathcal{L}_{reg} = \mathcal{L} + \frac{\lambda}{2} \|w\|_2^2$，使权重趋向较小值。
    \item \textbf{L1 正则化}（Lasso）：$\mathcal{L}_{reg} = \mathcal{L} + \lambda \|w\|_1$，产生稀疏权重，可用于特征选择。
    \item \textbf{Elastic Net}：L1 + L2 的组合。
\end{itemize}

\section{多分类扩展：Softmax 回归}
当类别数 $K > 2$ 时，逻辑回归可扩展为 \textbf{Softmax 回归}（多项逻辑回归）：
$$ P(y = k | x) = \frac{e^{w_k^\top x + b_k}}{\sum_{j=1}^K e^{w_j^\top x + b_j}} $$
\textbf{变量含义}：
\begin{itemize}
    \item $K$：总类别数。
    \item $w_k, b_k$：第 $k$ 类的权重向量和偏置。
    \item 分母是归一化常数，确保所有类别概率之和为 1。
\end{itemize}
对应的损失函数为\textbf{多类交叉熵}：
$$ \mathcal{L} = -\frac{1}{N} \sum_{i=1}^N \sum_{k=1}^K y_{ik} \log P(y_i = k | x_i) $$

\section{关键代码片段}
\begin{lstlisting}
import numpy as np

class LogisticRegression:
    def __init__(self, lr=0.01, n_iters=1000):
        self.lr = lr
        self.n_iters = n_iters
        self.w = None
        self.b = None

    def _sigmoid(self, z):
        return 1 / (1 + np.exp(-z))

    def fit(self, X, y):
        n_samples, n_features = X.shape
        self.w = np.zeros(n_features)
        self.b = 0
        for _ in range(self.n_iters):
            z = np.dot(X, self.w) + self.b
            y_hat = self._sigmoid(z)
            # 梯度 = (1/N) * X^T * (y_hat - y)
            dw = (1 / n_samples) * np.dot(X.T, (y_hat - y))
            db = (1 / n_samples) * np.sum(y_hat - y)
            self.w -= self.lr * dw
            self.b -= self.lr * db

    def predict(self, X):
        z = np.dot(X, self.w) + self.b
        return (self._sigmoid(z) >= 0.5).astype(int)
\end{lstlisting}

\section{逻辑回归 vs 其他模型}
\begin{table}[h]
\centering
\begin{tabularx}{\textwidth}{|l|X|}
\hline
\textbf{对比维度} & \textbf{分析} \\
\hline
\textbf{vs 线性回归} & 线性回归解决回归问题，逻辑回归解决分类问题。逻辑回归多了一个 Sigmoid 激活。 \\
\hline
\textbf{vs SVM} & SVM 关注最大化间隔（几何直觉），逻辑回归关注最大化似然（概率直觉）。SVM 不直接输出概率。 \\
\hline
\textbf{vs 决策树} & 决策树能处理非线性，但容易过拟合。逻辑回归是线性模型，更稳健但表达力有限。 \\
\hline
\textbf{vs 神经网络} & 单个神经元 + Sigmoid 激活 = 逻辑回归。神经网络是多层逻辑回归的堆叠。 \\
\hline
\end{tabularx}
\end{table}

\section{深度面试题}

\textbf{问：逻辑回归的损失函数为什么是凸的？}\\
答：BCE 损失函数的 Hessian 矩阵为 $H = \frac{1}{N} X^\top D X$，其中 $D = \text{diag}(\hat{y}_i(1-\hat{y}_i))$ 是一个对角矩阵且对角元素恒正。因此 $H$ 是半正定矩阵，损失函数是凸函数。

\textbf{问：逻辑回归能否处理多重共线性？}\\
答：逻辑回归对多重共线性敏感。当特征高度相关时，权重估计的方差会变大，模型不稳定。解决方法：1. L2 正则化；2. 主成分分析（PCA）降维；3. 特征选择。

\section{参考文献与拓展}
\begin{enumerate}
    \item \textbf{经典教材}: Bishop, C. M. (2006). \textit{Pattern Recognition and Machine Learning}, Chapter 4.
    \item \textbf{统计视角}: Hastie, T., et al. (2009). \textit{The Elements of Statistical Learning}, Chapter 4.3.
    \item \textbf{工程实践}: \texttt{sklearn.linear\_model.LogisticRegression} 官方文档。
\end{enumerate}

\end{document}
